% ============================================================
\setpartpreamble{\partintro
  {Grenzen der Laufzeit-Haken und weitere Erweiterungspunkte: Massenoperationen, Value
  Generation, Update-Pipeline, Complex Types, Modelle pro Laufzeitkontext und Migrations für
  Fortgeschrittene.}
  {
    \item Kapitel~\ref{ch:bulk-operations-and-the-limits}: Massenoperationen und die Grenzen
    von Interceptors
    \item Kapitel~\ref{ch:value-generation}: Value Generation
    \item Kapitel~\ref{ch:the-update-pipeline}: Die Update-Pipeline
    \item Kapitel~\ref{ch:complex-types-and-json}: Complex Types und JSON
    \item Kapitel~\ref{ch:models-per-runtime-context}: Modelle pro Laufzeitkontext
    \item Kapitel~\ref{ch:advanced-migrations}: Migrations für Fortgeschrittene
  }}
\part{Fortgeschrittene Erweiterungspunkte}\label{part:advanced}

% ------------------------------------------------------------
\chapter{Massenoperationen und die Grenzen von Interceptors}\label{ch:bulk-operations-and-the-limits}
\versionrisk{medium}

Ein SaveChanges-Interceptor sieht nur, was über \code{Save\allowbreak{}Changes} läuft.
\code{Execute\allowbreak{}Update} und \code{Execute\allowbreak{}Delete} schreiben direkt in
die Datenbank, ohne Change Tracker und ohne SaveChanges. Dasselbe gilt für rohes SQL. Ein
Plugin, das seine Garantien nur über SaveChanges durchsetzt, hat hier eine Lücke:

\begin{datatable}{H{1.28}H{0.88}H{0.88}H{0.96}}
  \textbf{Weg} & \textbf{SaveChanges-Interceptor} & \textbf{Connection-Interceptor} & \textbf{RLS} \\
  \tablehead
  \code{Save\allowbreak{}Changes} & ja & ja & ja \\
  \tablerule
  \code{Execute\allowbreak{}Update}, \code{Execute\allowbreak{}Delete} & \textbf{nein} & ja
  & ja \\
  \tablerule
  \code{From\allowbreak{}Sql}, \code{Sql\allowbreak{}Query}, \code{Execute\allowbreak{}Sql}
  & \textbf{nein} & ja & ja \\
  \tablerule
  ADO.NET oder psql als App-Rolle & \textbf{nein} & \textbf{nein} & ja (ohne Mandant: keine
  Zeilen) \\
\end{datatable}

Im Beispiel fängt RLS jeden dieser Wege ab, weil die Datenbank die Regeln selbst durchsetzt.
Ein Integrationstest beweist das, indem er den Interceptor gezielt umgeht:

% source: tests/Acme.EntityFrameworkCore.Tenancy.IntegrationTests/AdvancedTests.cs
\begin{sourcecode}[tests/\allowbreak{}.\allowbreak{}.\allowbreak{}.\allowbreak{}Integration\allowbreak{}Tests/\allowbreak{}Advanced\allowbreak{}Tests.\allowbreak{}cs (\allowbreak{}Auszug)]{csharp}
[Fact]
public async Task ExecuteUpdate_bypasses_interceptor_but_not_rls()
{
    // ... Bestellung für "tenant" anlegen ...
    await using var attacker = ContextFor(tenant);

    // Kein SaveChanges -> der SaveChanges-Interceptor läuft NICHT.
    // WITH CHECK der Policy verhindert das Umhängen trotzdem.
    var ex = await Assert.ThrowsAsync<PostgresException>(() =>
        attacker.Orders.ExecuteUpdateAsync(
            s => s.SetProperty(o => o.TenantId, TenantId.New())));

    Assert.Equal(PostgresErrorCodes.InsufficientPrivilege, ex.SqlState);
}
\end{sourcecode}

\begin{important}[Was ein Plugin dokumentieren muss]
  Plugins ohne Datenbank-Absicherung (Soft Delete, Audit-Felder, Mandantenfilter nur per
  Query-Filter) müssen dokumentieren, dass Massenoperationen und rohes SQL ihre Regeln
  umgehen. Wer das verhindern will, kann einen
  \code{IQuery\allowbreak{}Expression\allowbreak{}Interceptor} einsetzen:
  \code{Execute\allowbreak{}Update} und \code{Execute\allowbreak{}Delete} laufen durch die
  Query-Pipeline, der Interceptor sieht ihren Ausdrucksbaum und kann etwa ein
  \code{Set\allowbreak{}Property} auf geschützte Spalten ablehnen. Rohes SQL bleibt trotzdem
  offen.
\end{important}

\section{Beispiel: die Mandantenspalte bei Massen-Updates schützen}

Das Beispiel setzt genau diesen Interceptor um, als Option:
\code{Use\allowbreak{}Acme\allowbreak{}Tenancy(\allowbreak{}t => t.\allowbreak{}Protect\allowbreak{}Bulk\allowbreak{}Operations(\allowbreak{}))}
für PostgreSQL, \code{sql.\allowbreak{}Protect\allowbreak{}Bulk\allowbreak{}Operations(\allowbreak{})} für SQL
Server. EF Core~10 nimmt die Setter von \code{Execute\allowbreak{}Update} nicht mehr als
Ausdrucksbaum, sondern als \code{Action<\allowbreak{}Update\allowbreak{}Setters\allowbreak{}Builder<\allowbreak{}T>>}
entgegen. Im Baum, den der Interceptor sieht, stehen sie danach als Array von Tupeln:
\code{Execute\allowbreak{}Update(\allowbreak{}source,\allowbreak{} new[] \{ new Tuple(\allowbreak{}e => e.\allowbreak{}Tenant\allowbreak{}Id,\allowbreak{} wert),\allowbreak{} .\allowbreak{}.\allowbreak{}. \})}.
Das erste Element jedes Tupels ist das Property-Lambda, und genau das prüft der Visitor:

% source: src/Acme.EntityFrameworkCore.Tenancy/Interceptors/TenantQueryExpressionInterceptor.cs
\begin{sourcecode}[Interceptors/\allowbreak{}Tenant\allowbreak{}Query\allowbreak{}Expression\allowbreak{}Interceptor.\allowbreak{}cs (\allowbreak{}Kernpaket,\allowbreak{} gekürzt)]{csharp}
protected override Expression VisitMethodCall(MethodCallExpression node)
{
    if (node.Method.DeclaringType == typeof(EntityFrameworkQueryableExtensions)
        && node.Method.Name is nameof(EntityFrameworkQueryableExtensions.ExecuteUpdate)
            or nameof(EntityFrameworkQueryableExtensions.ExecuteUpdateAsync)
        && node.Arguments.Count > 1
        && node.Arguments[1] is NewArrayExpression setters)
    {
        foreach (var setter in setters.Expressions)
        {
            if (setter is NewExpression { Arguments: [var selector, ..] }
                && Unquote(selector) is LambdaExpression lambda)
                Check(lambda);
        }
    }

    return base.VisitMethodCall(node);
}

private void Check(LambdaExpression selector)
{
    var parameter = selector.Parameters[0];
    var name = selector.Body switch
    {
        // e => e.TenantId
        MemberExpression member when member.Expression == parameter => member.Member.Name,
        // e => EF.Property<TenantId>(e, "TenantId"), auch für Shadow-Properties
        MethodCallExpression
        {
            Method.Name: nameof(EF.Property),
            Arguments: [var target, ConstantExpression { Value: string property }]
        } when target == parameter => property,
        _ => null,
    };
    // ...
}
\end{sourcecode}

Ist die gesetzte Property die Mandanten-Property laut Laufzeit-Annotation, wirft der
Interceptor eine \code{Tenancy\allowbreak{}Exception}, noch bevor eine Verbindung geöffnet wird.
Er läuft nur beim Kompilieren, also einmal pro Abfrageform vor der Übersetzung; Ausführungen
aus dem Query-Cache kosten nichts. Andere Spalten und gewöhnliche Abfragen passieren unverändert. RLS bleibt die
eigentliche Grenze: Der Interceptor verhindert nur das Umhängen von Zeilen früh und mit klarer
Meldung, \code{Execute\allowbreak{}Delete} und rohes SQL prüft er nicht.

Eine Falle steckt in der Registrierung. \code{IQuery\allowbreak{}Expression\allowbreak{}Interceptor}
erbt wie \code{IMaterialization\allowbreak{}Interceptor} von
\code{ISingleton\allowbreak{}Interceptor}. Im internen Provider unter diesem Typ registriert,
ruft EF ihn aber nie auf: Query-Expression-Interceptors holt EF aus den
\code{IInterceptor}-Registrierungen und den Options, Materialisierungs-Interceptors auch aus
den \code{ISingleton\allowbreak{}Interceptor}-Registrierungen. Das Plugin registriert ihn deshalb
als \code{IInterceptor}; zustandslos, wie er ist, schadet die Instanz pro Context nicht.
\code{Bulk\allowbreak{}Operation\allowbreak{}Tests} hält beide Registrierungen fest, die
Integrationstests zeigen ein erlaubtes und ein abgelehntes Update gegen PostgreSQL.

% ------------------------------------------------------------
\chapter{Value Generation}\label{ch:value-generation}
\versionrisk{stable}

EF erzeugt Schlüsselwerte entweder in der Datenbank (Identity, Sequenzen) oder clientseitig
über einen \code{Value\allowbreak{}Generator}. Für \code{Guid}-Schlüssel macht EF das von
selbst, für eigene Typen wie \code{Tenant\allowbreak{}Id} nicht. Ein Plugin schließt diese
Lücke mit einem Generator, einer Factory und einer Convention:

% source: src/Acme.EntityFrameworkCore.Tenancy/ValueGeneration/TenantIdValueGenerator.cs
\begin{sourcecode}[Value\allowbreak{}Generation/\allowbreak{}Tenant\allowbreak{}Id\allowbreak{}Value\allowbreak{}Generator.\allowbreak{}cs (\allowbreak{}Kernpaket)]{csharp}
public class TenantIdValueGenerator : ValueGenerator<TenantId>
{
    // Endgültige Werte, keine temporären Platzhalter bis zum Speichern
    public override bool GeneratesTemporaryValues => false;

    public override TenantId Next(EntityEntry entry) => TenantId.New();   // UUIDv7
}
\end{sourcecode}

% source: src/Acme.EntityFrameworkCore.Tenancy/ValueGeneration/TenantIdValueGeneratorFactory.cs
\begin{sourcecode}[Value\allowbreak{}Generation/\allowbreak{}Tenant\allowbreak{}Id\allowbreak{}Value\allowbreak{}Generator\allowbreak{}Factory.\allowbreak{}cs (\allowbreak{}Kernpaket)]{csharp}
public class TenantIdValueGeneratorFactory : ValueGeneratorFactory
{
    // Zustandsloser Generator: EF darf ihn cachen und zwischen Contexts teilen
    public override ValueGenerator Create(IProperty property, ITypeBase typeBase)
        => new TenantIdValueGenerator();
}
\end{sourcecode}

% source: src/Acme.EntityFrameworkCore.Tenancy/Conventions/TenantIdKeyValueGenerationConvention.cs
\begin{sourcecode}[Conventions/\allowbreak{}Tenant\allowbreak{}Id\allowbreak{}Key\allowbreak{}Value\allowbreak{}Generation\allowbreak{}Convention.\allowbreak{}cs (\allowbreak{}Kernpaket)]{csharp}
public virtual void ProcessModelFinalizing(
    IConventionModelBuilder modelBuilder,
    IConventionContext<IConventionModelBuilder> context)
{
    foreach (var entityType in modelBuilder.Metadata.GetEntityTypes())
    {
        if (entityType.BaseType is not null
            || entityType.FindPrimaryKey() is not { Properties: [var property] }
            || property.ClrType != typeof(TenantId))
            continue;

        // Convention-Quelle: ein explizites ValueGeneratedNever() gewinnt (-> null)
        if (property.Builder.ValueGenerated(ValueGenerated.OnAdd) is null)
            continue;

        property.Builder.HasValueGeneratorFactory(typeof(TenantIdValueGeneratorFactory));
    }
}
\end{sourcecode}

\begin{important}[Factory-Typ statt Generator]
  Einen Value Generator kann man auf mehreren Wegen konfigurieren:
  \code{Has\allowbreak{}Value\allowbreak{}Generator} mit einem Typ oder einer
  Factory-Funktion, oder
  \code{Has\allowbreak{}Value\allowbreak{}Generator\allowbreak{}Factory} mit einem
  Factory-Typ. Der Generator für Compiled Models in EF Core 10 akzeptiert nur die letzte
  Variante und bricht sonst mit „The property ... has a value generator configured. Use
  'HasValueGeneratorFactory' to configure the value generator factory type.“ ab. Plugins
  sollten deshalb
  \code{Has\allowbreak{}Value\allowbreak{}Generator\allowbreak{}Factory(\allowbreak{}Type)}
  verwenden und der Factory einen öffentlichen parameterlosen Konstruktor geben.
\end{important}

Die Convention liegt im Kernpaket, weil sie nichts mit PostgreSQL zu tun hat; das
Provider-Paket hängt sie über sein Convention-Set-Plugin ein. Der Compiled-Model-Test prüft,
dass die Factory im erzeugten Code auftaucht.

% ------------------------------------------------------------
\chapter{Die Update-Pipeline}\label{ch:the-update-pipeline}
\versionrisk{volatile}

Kapitel~\ref{ch:the-query-pipeline} beschreibt den Weg von LINQ zu SQL. Für das Schreiben
gibt es eine eigene Pipeline, die Plugins für Massenoperationen, datenbankseitige
Soft-Deletes oder zurückgelieferte Werte erweitern:

\begin{flowdiagram}
  \flowstep{ChangeTracker}{Ermittelt hinzugefügte, geänderte und gelöschte Entitäten}
  \flowstep{IModificationCommandFactory}{Ein ModificationCommand pro Zeile und Tabelle}
  \flowstep{ICommandBatchPreparer}{Sortiert nach Abhängigkeiten (Fremdschlüssel) und bildet
  Batches}
  \flowstep{IModificationCommandBatchFactory}{Erzeugt die Batches (Größe,
  Provider-Besonderheiten)}
  \flowstep{IUpdateSqlGenerator}{Schreibt INSERT, UPDATE, DELETE, ggf. mit RETURNING}
  \flowstep{Ausführung}{Ergebnisse (Identity, Concurrency-Token) zurück in die Entitäten}
\end{flowdiagram}

\begin{datatable}{H{1.64}H{0.68}H{0.68}}
  \textbf{Dienst} & \textbf{Lebensdauer} & \textbf{Art} \\
  \tablehead
  \code{IModification\allowbreak{}Command\allowbreak{}Factory} & Singleton & einzeln \\
  \tablerule
  \code{ICommand\allowbreak{}Batch\allowbreak{}Preparer} & Scoped & einzeln \\
  \tablerule
  \code{IModification\allowbreak{}Command\allowbreak{}Batch\allowbreak{}Factory} & Scoped &
  einzeln \\
  \tablerule
  \code{IUpdate\allowbreak{}Sql\allowbreak{}Generator} & Singleton & einzeln \\
\end{datatable}

Alle Dienste sind Einzeldienste des Providers, mit den Konsequenzen aus
Kapitel~\ref{ch:registering-services}. Bevor man sie ersetzt, lohnt die Prüfung der
Alternativen: Ein SaveChanges-Interceptor kann Zustände umschreiben (aus Löschen wird Ändern
beim Soft Delete), ein Command-Interceptor kann SQL ergänzen, und für Massenänderungen gibt
es \code{Execute\allowbreak{}Update}. Ein eigener
\code{IUpdate\allowbreak{}Sql\allowbreak{}Generator} ist erst dann gerechtfertigt, wenn das
SQL selbst eine andere Form braucht, etwa \code{INSERT
.\allowbreak{}.\allowbreak{}.\allowbreak{} ON CONFLICT} für Upserts.

\section{Beispiel: die Datenbankuhr beim UPDATE}

Eine Spalte \code{Updated\allowbreak{}At} soll den Zeitpunkt der letzten Änderung tragen, und
darüber soll die Uhr der Datenbank entscheiden, nicht die des Anwendungsservers, der die Zeile
gerade schreibt. Für \code{INSERT} hat EF das Mittel:
\code{Has\allowbreak{}Default\allowbreak{}Value\allowbreak{}Sql(\allowbreak{}\textquotedbl{}now()\textquotedbl{})}. Für
\code{UPDATE} gibt es keines.
\code{Value\allowbreak{}Generated\allowbreak{}On\allowbreak{}Add\allowbreak{}Or\allowbreak{}Update(\allowbreak{})}
sagt EF nur, dass die Datenbank den Wert erzeugt; EF lässt die Spalte im \code{SET} weg, liest
sie mit \code{RETURNING} zurück und erwartet, dass ein Trigger sie gesetzt hat. Die einfacheren
Mittel erreichen das Ziel nicht:

\begin{itemize}
  \item Ein \textbf{SaveChanges-Interceptor} kann nur einen Wert aus der Anwendung setzen: die
    Uhr des Anwendungsservers, nicht \code{now()} der Transaktion.
  \item Ein \textbf{Command-Interceptor} müsste fertigen SQL-Text bearbeiten, mit
    Parameternamen und \code{RETURNING}-Klausel, in jedem Batch: fragil.
  \item Ein \textbf{Trigger} funktioniert, ist aber ein Schemaobjekt außerhalb des EF-Modells,
    das jede Migration mitführen muss.
\end{itemize}

Es fehlt eine andere Form des \code{UPDATE} selbst, und die schreibt nur der
\code{IUpdate\allowbreak{}Sql\allowbreak{}Generator}. Das Beispiel ist ein Lehrbeispiel im
Testprojekt, nicht Teil der Tenancy-Pakete. Eine Fluent-API-Methode markiert die Property, der
abgeleitete Generator ergänzt die Zuweisung:

% source: tests/Acme.EntityFrameworkCore.Tenancy.IntegrationTests/UpdatePipeline/DatabaseClockUpdateSqlGenerator.cs
\begin{sourcecode}[tests/\allowbreak{}.\allowbreak{}.\allowbreak{}.\allowbreak{}/\allowbreak{}Update\allowbreak{}Pipeline/\allowbreak{}Database\allowbreak{}Clock\allowbreak{}Update\allowbreak{}Sql\allowbreak{}Generator.\allowbreak{}cs (\allowbreak{}gekürzt)]{csharp}
public static PropertyBuilder<T> HasDatabaseClockOnUpdate<T>(this PropertyBuilder<T> builder)
{
    builder.Metadata.SetAnnotation(Annotation, true);
    return builder.HasDefaultValueSql("now()").ValueGeneratedOnAddOrUpdate();
}

public class DatabaseClockUpdateSqlGenerator : NpgsqlUpdateSqlGenerator   // intern (EF1001)
{
    protected override void AppendUpdateCommand(
        StringBuilder commandStringBuilder,
        string name,
        string? schema,
        IReadOnlyList<IColumnModification> writeOperations,
        IReadOnlyList<IColumnModification> readOperations,
        IReadOnlyList<IColumnModification> conditionOperations,
        bool appendReturningOneClause = false)
    {
        AppendUpdateCommandHeader(commandStringBuilder, name, schema, writeOperations);

        // EFs Reihenfolge, mit den zusätzlichen Zuweisungen nach SET
        foreach (var column in readOperations.Where(IsDatabaseClock))
        {
            commandStringBuilder.Append(", ");
            SqlGenerationHelper.DelimitIdentifier(commandStringBuilder, column.ColumnName);
            commandStringBuilder.Append(" = now()");
        }

        AppendWhereClause(commandStringBuilder, conditionOperations);
        AppendReturningClause(commandStringBuilder, readOperations,
            appendReturningOneClause ? "1" : null);
        commandStringBuilder.AppendLine(SqlGenerationHelper.StatementTerminator);
    }
}
\end{sourcecode}

Die Überschreibung ändert so wenig wie möglich. Eine beim Ändern generierte Spalte ist für EF
eine Leseoperation: Sie steht im \code{RETURNING}, nicht im \code{SET}. Die Methode schreibt den
Kopf von EF, hängt für jede markierte Lesespalte \code{\textquotedbl{}Spalte\textquotedbl{} = now()} an und überlässt die
\code{WHERE}-Klausel mit Schlüssel und Concurrency-Token sowie die \code{RETURNING}-Klausel EF.
Deshalb funktionieren der zurückgelesene Wert und die optimistische Konkurrenzprüfung weiter.
Mit \code{xmin} als Row Version sendet SaveChanges:

\begin{sourcecode}[SaveChanges nach einer Änderung von Title]{sql}
UPDATE "Documents" SET "Title" = @p0, "UpdatedAt" = now()
WHERE "Id" = @p1 AND xmin = @p2
RETURNING "UpdatedAt", xmin;
\end{sourcecode}

\code{IUpdate\allowbreak{}Sql\allowbreak{}Generator} ist ein Singleton (Tabelle oben): Der
Generator darf keinen Zustand pro Befehl halten; alles, was er braucht, kommt aus den
Parametern. Er ersetzt einen Provider-Dienst mit \code{Replace\allowbreak{}Service}, hier neben
dem Tenancy-Plugin, das diesen Dienst nicht berührt:

\begin{sourcecode}[Registrierung]{csharp}
options.UseNpgsql(connectionString)
    .UseAcmeTenancy()
    .ReplaceService<IUpdateSqlGenerator, DatabaseClockUpdateSqlGenerator>();

modelBuilder.Entity<Document>(b =>
{
    b.Property(d => d.UpdatedAt).HasDatabaseClockOnUpdate();
    b.Property(d => d.Version).IsRowVersion();            // Npgsql: die Systemspalte xmin
});
\end{sourcecode}

Wie weit das Beispiel trägt, entscheidet die Basisklasse. Sie leitet von
\code{Npgsql\allowbreak{}Update\allowbreak{}Sql\allowbreak{}Generator} ab, einer internen
Npgsql-Klasse (EF1001): Mit EFs Basisklasse gingen Npgsqls eigene SQL-Formen verloren, etwa
Teilaktualisierungen von JSON-Spalten. Daher rührt die Einstufung dieses Kapitels. Die Grenzen:

\begin{itemize}
  \item \textbf{Einer pro Context.} Es gibt genau einen Update-SQL-Generator. Ein zweites
    Plugin, das ihn ersetzt, gewinnt oder verliert je nach Reihenfolge der Registrierung;
    kombinieren lassen sich beide nur, indem einer vom anderen ableitet.
  \item \textbf{Nur SaveChanges.} \code{Execute\allowbreak{}Update} läuft nicht durch diese
    Pipeline (Kapitel~\ref{ch:bulk-operations-and-the-limits}); es muss die Spalte selbst
    setzen. Auf Stored Procedures abgebildete Entitäten ebenfalls nicht.
  \item \textbf{Providerspezifisch.} Die Überschreibung setzt auf \code{RETURNING} von
    PostgreSQL; SQL Server schreibt \code{OUTPUT} und braucht einen eigenen Generator.
\end{itemize}

\code{Update\allowbreak{}Pipeline\allowbreak{}Tests} prüft das Beispiel gegen PostgreSQL: Das
aufgezeichnete \code{UPDATE} enthält die Zuweisung, das \code{WHERE} mit \code{xmin} und das
\code{RETURNING}; die Entität trägt danach genau den Wert der Datenbank; und ein zweiter
Schreiber mit veralteter Version erhält eine
\code{Db\allowbreak{}Update\allowbreak{}Concurrency\allowbreak{}Exception} und ändert nichts.
Ohne \code{Replace\allowbreak{}Service} scheitert der SQL-Test.

% ------------------------------------------------------------
\chapter{Complex Types und JSON}\label{ch:complex-types-and-json}
\versionrisk{medium}

\section{Complex Types}

Complex Types (seit EF Core 8) sind Werttypen ohne eigene Identität, deren Properties in der
Tabelle des Besitzers landen. Für Plugins heißt das: Conventions, die über
\code{entity\allowbreak{}Type.\allowbreak{}Get\allowbreak{}Properties(\allowbreak{})}
laufen, sehen deren Properties \textbf{nicht}. Sie hängen an
\code{IComplex\allowbreak{}Property} und müssen eigens durchlaufen werden:

\begin{sourcecode}[Muster für Conventions (\allowbreak{}Skizze)]{csharp}
foreach (var complexProperty in entityType.GetComplexProperties())
{
    foreach (var property in complexProperty.ComplexType.GetProperties())
    {
        // Spaltennamen sind präfixiert (z. B. "Address_City"); immer über
        // property.GetColumnName(table) ermitteln, nie selbst zusammensetzen
    }
}
\end{sourcecode}

Fluent-API-Methoden für Properties sollten zusätzlich auf
\code{Complex\allowbreak{}Type\allowbreak{}Property\allowbreak{}Builder} angeboten werden,
wenn sie auch dort sinnvoll sind. Das Beispiel-Plugin braucht das nicht: Isolation und
Wertgenerierung betreffen nur Entitäten.

\section{JSON-Spalten}

Owned Types und (seit EF Core 10) Complex Types lassen sich mit
\code{To\allowbreak{}Json(\allowbreak{})} als eine JSON-Spalte speichern. Werte darin liest
und schreibt EF über den
\code{Json\allowbreak{}Value\allowbreak{}Reader\allowbreak{}Writer} des Type Mappings
(Kapitel~\ref{ch:type-mappings-and-value-converters}). Ein Value Converter allein genügt nur,
wenn EF den Reader/Writer des Providertyps darum legen kann
(\code{Json\allowbreak{}Converted\allowbreak{}Value\allowbreak{}Reader\allowbreak{}Writer});
ein eigenes Mapping ohne beides lässt den Modellbau scheitern.
Unit- und Integrationstest des Beispiels prüfen beide Seiten:

% source: tests/Acme.EntityFrameworkCore.Tenancy.Tests/JsonTests.cs
\begin{sourcecode}[tests/\allowbreak{}.\allowbreak{}.\allowbreak{}.\allowbreak{}/\allowbreak{}Json\allowbreak{}Tests.\allowbreak{}cs]{csharp}
using var ctx = Contexts.Create(model: m =>
    m.Entity<Shipment>().OwnsOne(s => s.Details, d => d.ToJson()));

var property = ctx.DesignModel()
    .FindEntityType(typeof(ShipmentDetails))!
    .FindProperty(nameof(ShipmentDetails.ApprovedBy))!;

Assert.Same(TenantIdJsonReaderWriter.Instance,
    property.GetTypeMapping().JsonValueReaderWriter);
\end{sourcecode}

Davon getrennt ist JSON außerhalb der Datenbank. System.Text.Json kennt den
\code{Json\allowbreak{}Value\allowbreak{}Reader\allowbreak{}Writer} von EF nicht und würde
\code{TenantId} als Objekt mit \code{Value} und \code{IsEmpty} schreiben. Ein
\code{[Json\allowbreak{}Converter(\allowbreak{}typeof(\allowbreak{}Tenant\allowbreak{}Id\allowbreak{}Json\allowbreak{}Converter))]}
am Typ im Kernpaket macht daraus eine schlichte GUID, auch als Dictionary-Schlüssel, ohne dass
eine Web-API etwas registrieren muss; ungültige Werte führen zu einer
\code{Json\allowbreak{}Exception}. Beide Wege brauchen ihren eigenen Baustein, weil EF und
System.Text.Json unabhängig voneinander serialisieren
(\code{Tenant\allowbreak{}Id\allowbreak{}Json\allowbreak{}Converter\allowbreak{}Tests}).

\begin{note}
  PostgreSQL-Arrays (\code{List<\allowbreak{}Tenant\allowbreak{}Id>} als \code{uuid[]}) sind
  kein JSON: Npgsql bildet primitive Collections nativ als Arrays ab. Ob ein Provider
  Collections mit konvertierten Elementtypen unterstützt, ist providerspezifisch und gehört
  in einen eigenen Test, bevor man es dokumentiert.
\end{note}

% ------------------------------------------------------------
\chapter{Modelle pro Laufzeitkontext}\label{ch:models-per-runtime-context}
\versionrisk{stable}

EF baut das Modell einmal pro internem Service Provider und Context-Typ und cacht es
(Kapitel~\ref{ch:ef-core-s-architecture-from}). Soll das Modell von Laufzeitdaten abhängen,
etwa einem Schema pro Mandant, muss der Cache-Schlüssel diese Daten enthalten. Dafür gibt es
\code{IModel\allowbreak{}Cache\allowbreak{}Key\allowbreak{}Factory}:

\begin{sourcecode}[tests/\allowbreak{}.\allowbreak{}.\allowbreak{}.\allowbreak{}/\allowbreak{}Schema\allowbreak{}Per\allowbreak{}Tenant\allowbreak{}Tests.\allowbreak{}cs (\allowbreak{}gekürzt, vereinfacht)]{csharp}
public sealed class SchemaModelCacheKeyFactory : IModelCacheKeyFactory
{
    public object Create(DbContext context, bool designTime)
        => context is ISchemaScoped scoped
            ? (context.GetType(), scoped.Schema, designTime)
            : (object)(context.GetType(), designTime);
}

public sealed class SchemaContext : DbContext, ISchemaScoped
{
    // Konstruktor: (DbContextOptions<SchemaContext> options, string schema)
    public string Schema { get; }

    protected override void OnModelCreating(ModelBuilder modelBuilder)
        => modelBuilder.HasDefaultSchema(Schema);
}

// Registrierung
options.UseNpgsql(connectionString)
       .ReplaceService<IModelCacheKeyFactory, SchemaModelCacheKeyFactory>();
\end{sourcecode}

Der Test \code{Schema\allowbreak{}Per\allowbreak{}Tenant\allowbreak{}Tests} prüft beide
Seiten des Cache-Schlüssels: Zwei Schemas ergeben zwei Modelle mit jeweils eigenem
Standardschema, und dasselbe Schema verwendet das gecachte Modell wieder, statt es neu zu
bauen.

Das ist die Grundlage für „Schema pro Mandant“, eine Alternative zu RLS. Die drei üblichen
Ansätze im Vergleich:

\begin{datatable}{H{0.83}H{1.02}H{1.15}}
  \textbf{Ansatz} & \textbf{Stärken} & \textbf{Schwächen} \\
  \tablehead
  Gemeinsame Tabellen + RLS (Beispiel) & Ein Modell, ein Schema, einfache Migrationen;
  Trennung erzwingt die Datenbank & Alle Mandanten teilen Tabellen und Indizes;
  Fehlkonfiguration der Rollen hebelt die Trennung aus \\
  \tablerule
  Schema pro Mandant & Starke Trennung, getrennte Backups pro Schema möglich & Ein Modell
  pro Mandant im Cache (Speicher!), Migrationen für jedes Schema \\
  \tablerule
  Datenbank pro Mandant & Maximale Trennung, unabhängige Skalierung & Ein Modell genügt,
  aber Verbindungen, Migrationen und Betrieb pro Datenbank \\
\end{datatable}

\begin{important}[Modell-Cache und Speicher]
  Jeder Schlüsselwert erzeugt ein eigenes Modell im Speicher. Bei Hunderten Mandanten mit
  großen Modellen wird das teuer. Außerdem dürfen Werte, die das Modell verändern, nie aus
  ungeprüften Eingaben stammen.
\end{important}

% ------------------------------------------------------------
\chapter{Migrations für Fortgeschrittene}\label{ch:advanced-migrations}
\versionrisk{stable}

\section{Reihenfolge von Operationen}

Generierte und handgeschriebene Operationen laufen in der Reihenfolge, in der sie in
\code{Up(\allowbreak{})} stehen. Das ist eine echte Fehlerquelle. Die Policy
\code{support\_\allowbreak{}read} des Beispiels gilt \code{TO shop\_\allowbreak{}support}.
EF erzeugt sie direkt nach \code{CREATE TABLE}. Wird die Rolle erst am Ende von
\code{Up(\allowbreak{})} angelegt, scheitert die Migration.

\begin{sourcecode}[Migrations/\allowbreak{}20261021\_\allowbreak{}Initial.\allowbreak{}cs (\allowbreak{}handergänzt, vereinfacht)]{csharp}
protected override void Up(MigrationBuilder migrationBuilder)
{
    // 1. Was Policies referenzieren: VOR den generierten Operationen
    migrationBuilder.CreateApplicationRole("shop_support");

    // 2. Von EF generiert: Tabellen, RLS, Policies (TO shop_support)
    migrationBuilder.CreateTable(name: "Customers", /* ... */);

    // 3. Was Tabellen voraussetzt: NACH den generierten Operationen
    migrationBuilder.CreateApplicationRole("shop_app", canLogin: true);
    migrationBuilder.GrantTablePrivileges("shop_app");   // auch künftige Tabellen, Sequenzen
}
\end{sourcecode}

\section{Idempotente Skripte}

\code{dotnet ef migrations script -{}-idempotent} erzeugt ein Skript, das gegen jeden
Datenbankstand laufen kann. Npgsql verpackt dazu jede Operation in einen Block \code{DO
\$EF\$ .\allowbreak{}.\allowbreak{}.\allowbreak{} \$EF\$}, der die Historientabelle prüft.
Eigene Operationen, die selbst einen \code{DO \$\$}-Block erzeugen, stecken dann darin. Das
funktioniert, weil die Dollar-Quotes verschiedene Tags haben. Solche Annahmen gehören in
einen Test:

% source: tests/Acme.EntityFrameworkCore.Tenancy.IntegrationTests/AdvancedTests.cs
\begin{sourcecode}[tests/\allowbreak{}.\allowbreak{}.\allowbreak{}.\allowbreak{}Integration\allowbreak{}Tests/\allowbreak{}Advanced\allowbreak{}Tests.\allowbreak{}cs (\allowbreak{}Auszug)]{csharp}
var script = db.GetService<IMigrator>()
    .GenerateScript(options: MigrationsSqlGenerationOptions.Idempotent);

await db.Database.ExecuteSqlRawAsync(script);   // verschachteltes DO gültig?
await db.Database.ExecuteSqlRawAsync(script);   // zweiter Lauf wird übersprungen
\end{sourcecode}

\section{Operationen außerhalb von Transaktionen}

Manche Befehle dürfen nicht in einer Transaktion laufen, in PostgreSQL etwa \code{CREATE
INDEX CONCURRENTLY}. Eigene Operationen mit solchem SQL beenden ihren Befehl mit
\code{End\allowbreak{}Statement(\allowbreak{}builder,\allowbreak{}
suppress\allowbreak{}Transaction: true)}. EF führt ihn dann außerhalb der
Migrationstransaktion aus. Der Preis: Schlägt die Migration danach fehl, wird dieser Befehl
nicht zurückgerollt; er sollte deshalb idempotent sein.

\section{Migration Bundles}

\code{dotnet ef migrations bundle} erzeugt eine ausführbare Datei, die Migrationen ohne SDK
anwendet. Sie enthält das Startprojekt samt Plugin-Assemblies; eigene Operationen und der
SQL-Generator des Plugins wirken also auch dort. Wichtig ist, dass die Design-Time-Factory
oder der Host dieselben Plugin-Optionen verwendet wie die Anwendung, sonst entsteht SQL für
ein anderes Modell. Wie das Bundle mit getrennten Rollen in eine Pipeline gehört, zeigt
Seite~\pageref{sec:deployment}.
